\documentclass[12pt]{article}
\usepackage[utf8]{inputenc}
\usepackage[T1]{fontenc}
\usepackage[margin=1in]{geometry}
\usepackage{setspace}
\usepackage{xcolor}
\usepackage[colorlinks=true, urlcolor=blue!60!black, linkcolor=blue!60!black]{hyperref}
\setlength{\parindent}{0em}
\setlength{\parskip}{0.6em}
\renewcommand{\baselinestretch}{1.0}
\begin{document}
\begin{center}
{\LARGE \textit{Brown v. Board of Education (Brown II)}, 349 U.S. 294 (1955)\par}
\vspace{0.5em}
{\large Case Brief --- Two-Page Edition\par}
\end{center}

\section{Memory Jogger}
Local school boards, supervised by federal courts sitting in equity, must desegregate and admit Black students on a racially nondiscriminatory basis ``with all deliberate speed'' (\href{https://supreme.justia.com/cases/federal/us/349/294/}{Justia}).

\section{Detailed Case Facts}
\textit{Brown I} (1954) held school segregation unconstitutional but set reargument on relief. The five consolidated cases came from Kansas, South Carolina, Virginia, the District of Columbia, and Delaware. Some of these communities had already begun desegregating; South Carolina and Virginia had done nothing. Southern leaders were mobilizing massive resistance while NAACP counsel urged firm deadlines.

\section{Procedural History}
The federal courts below had denied relief; only Delaware's courts had ordered immediate admission. After reargument in April 1955, the Court reversed and remanded the four federal cases and affirmed the Delaware judgment.

\section{Judicial Votes}
Unanimous, 9--0. Chief Justice Warren wrote for the Court; no Justice wrote separately.

\section{Holding}
All laws requiring or permitting segregation must yield to \textit{Brown I}. School authorities bear primary responsibility for desegregation plans; trial courts judge their good faith, require a ``prompt and reasonable start,'' place the burden of justifying delay on defendants, retain jurisdiction, and order admission ``with all deliberate speed.''

\section{Analysis of Opinions}
\textbf{Majority:} Warren treated the case as a question of remedy and institutional competence. Trial courts, close to local conditions, would judge whether boards acted in good faith. Equity would balance the plaintiffs' personal interest in prompt admission against the public interest in an orderly transition, but constitutional principles could not yield to disagreement. Once a prompt start was made, defendants could obtain more time only by proving it necessary, considering such obstacles as the school plant, transportation, personnel, and attendance areas. No Justice wrote separately.

\section{Examples: Future Applications}
\textbf{Same-Side:} a board that promptly redraws bus routes, reassigns teachers, and admits Black students next year has made a valid start. \textbf{Opposite-Side:} a three-year postponement sought because of public hostility must be denied. \textbf{Fence-Sitter:} a twelve-year grade-a-year plan starts at once but leaves students in upper grades without any desegregated schooling.

\section{Critique}
The Court entrusted desegregation to the boards that had defended segregation, and its oxymoronic formula invited delay until \textit{Alexander v. Holmes County} (1969) repudiated it; progressive scholars call it a tragic compromise. Originalist critics add that plaintiffs deserved an immediate remedy and that the decision launched open-ended structural injunctions.

\section{Key Quotations}
\begin{quote}
``Traditionally, equity has been characterized by a practical flexibility in shaping its remedies and by a facility for adjusting and reconciling public and private needs.''
\end{quote}
\noindent --- Chief Justice Warren, for the Court, 349 U.S. at 300 (footnotes omitted).

\end{document}
